\documentclass[11pt]{article}

\usepackage[T1]{fontenc}
\usepackage[utf8]{inputenc}
\usepackage{amsmath,amssymb}
\usepackage{booktabs}
\usepackage{geometry}
\usepackage[scaled]{helvet}
\usepackage{hyperref}

\geometry{margin=1in}
\hypersetup{hidelinks}
\renewcommand{\familydefault}{\sfdefault}
\setlength{\parindent}{0pt}

\title{Reproducibility of the Flagship Sweep}
\date{}

\begin{document}

\raggedright

\maketitle

\section{Summary}

We evaluated whether the completed sweep in \texttt{runs/flagship/} reproduces Figure
2 of \emph{Slow Transition to Low-Dimensional Chaos in Heavy-Tailed Recurrent Neural
Networks}. The answer depends on which Lyapunov analysis is used:

\begin{enumerate}
\item \textbf{The published numerical transition points are closely reproduced when we
emulate the original study's historical analysis errors.} This analysis
propagated tangent vectors with a transposed Jacobian, $J^\top Q$, and
evaluated the Jacobian one state transition too late. It reproduces 7 of the
paper's 9 transition grid points exactly; the other 2 differ by one adjacent
gain value.
\item \textbf{The results change when the Lyapunov analysis is corrected.} The canonical
update, $JQ$, with the Jacobian evaluated at the state that generates the
transition produces substantially earlier transitions in the most
heavy-tailed conditions. It exactly matches only 5 of the 9 reported grid
points and accumulates 11 grid steps of absolute transition error, compared
with 2 for the historical analysis.
\end{enumerate}

Thus, the paper's qualitative conclusion---a slower, broader transition to chaos
for heavier-tailed networks---remains visible under the corrected analysis. The
paper's reported quantitative transition locations, however, appear to reflect
the historical, erroneous Lyapunov calculation.

\section{Target experiment}

The comparison targets Figure 2 of the archived \href{original/paper.pdf}{paper},
which reports the maximum Lyapunov exponent as a function of recurrent gain for
autonomous networks with:

\begin{itemize}
\item $N\in\{1000,3000,10000\}$;
\item $\alpha\in\{1.0,1.5,2.0\}$;
\item 50 logarithmically spaced gains from $0.01$ to $10$;
\item 10 random trials per condition;
\item 2,900 warmup steps followed by 100 Lyapunov accumulation steps; and
\item the leading 100 Lyapunov exponents.
\end{itemize}

These settings are reproduced by
\href{configs/flagship.yaml}{\texttt{configs/flagship.yaml}}. The completed run contains all
4,500 expected condition/trial artifacts. Every recorded metric is finite, no
QR stretch floors were used, and the 52 non-CUDA tests pass.

The transition estimator follows the paper's plotting code: for each
$(N,\alpha)$, trial MLEs are averaged at every sampled gain, and $g^*$ is the
first gain where the mean MLE changes from negative to nonnegative. Because the
gain grid is logarithmic, adjacent values differ by a factor of approximately
1.151. A one-grid-step disagreement therefore corresponds to about a 15\%
difference in the reported discrete threshold.

\section{Correct and historical Lyapunov analyses}

For the autonomous recurrence

\[
h_{t+1}=\tanh(W h_t),
\]

the Jacobian for the transition from $h_t$ to $h_{t+1}$ is

\[
J_t=D_tW,
\qquad
D_t=\operatorname{diag}\!\left(\operatorname{sech}^2(W h_t)\right).
\]

The canonical forward tangent update is

\[
Q_{t+1}R_{t+1}=J_tQ_t.
\]

This is the algorithm stated in the paper and implemented by the corrected
code in \href{simulation.py\#L134-L141}{\texttt{simulation.py}}. The state update and its
Jacobian use the same preactivation.

The archived study implementation historically made two changes to this
calculation:

\begin{enumerate}
\item \textbf{Transposed tangent propagation.} A comment in
\href{original/LypAlgo.py\#L9-L12}{\texttt{original/LypAlgo.py}} records that the code used
$J^\top Q$ before July 8, 2026.
\item \textbf{One-step-late Jacobian indexing.} The loop first advances \texttt{ht} to
$h_{t+1}$ and then passes that updated state to \texttt{rnn\_jac}
(\href{original/LypAlgo.py\#L80-L84}{update},
\href{original/LypAlgo.py\#L96-L100}{Jacobian call}). \texttt{rnn\_jac} subsequently
recomputes $W h_{t+1}$ (\href{original/LypAlgo.py\#L138-L162}{construction}), so
the tangent update uses $J_{t+1}$ rather than $J_t$.
\end{enumerate}

The historical update was therefore effectively

\[
Q_{t+1}R_{t+1}=J_{t+1}^{\top}Q_t.
\]

This is not equivalent to the canonical calculation for a time-varying
Jacobian sequence. In particular,

\[
J_T^\top\cdots J_1^\top=(J_1\cdots J_T)^\top,
\]

which does not preserve the chronological product
$J_T\cdots J_1$. The difference is especially consequential for finite-time
Lyapunov estimates in nonlinear, strongly non-normal heavy-tailed networks.

\section{Finding 1: the historical analysis closely reproduces Figure 2}

We regenerated the weights from the saved condition seeds and re-evaluated the
saved trajectories using the historical calculation. The re-analysis used all
100 tangent directions, 100 float32 QR steps, the original $10^{-12}$ QR
floor, the maximum of the resulting 100 exponents per trial, and the mean over
10 trials. This matches the analysis dimensions used for Figure 2 while
holding the networks and trajectories fixed across analysis variants.

\begin{center}
\begin{tabular}{rrrrr}
\toprule
$N$ & $\alpha$ & Paper $g^*$ & Historical $J_{t+1}^\top Q$ & Correct $J_tQ$ \\
\midrule
1,000 & 1.0 & 0.791 & 0.910 & 0.391 \\
1,000 & 1.5 & 0.791 & 0.687 & 0.596 \\
1,000 & 2.0 & 0.910 & 0.910 & 0.910 \\
3,000 & 1.0 & 0.450 & 0.450 & 0.295 \\
3,000 & 1.5 & 0.596 & 0.596 & 0.596 \\
3,000 & 2.0 & 0.910 & 0.910 & 0.910 \\
10,000 & 1.0 & 0.295 & 0.295 & 0.256 \\
10,000 & 1.5 & 0.518 & 0.518 & 0.518 \\
10,000 & 2.0 & 0.791 & 0.791 & 0.791 \\
\bottomrule
\end{tabular}
\end{center}

The corresponding transition metrics are:

\begin{center}
\begin{tabular}{lrr}
\toprule
Analysis & Exact paper grid matches & Total absolute grid-step error \\
\midrule
Correct $J_tQ$ & 5/9 & 11 \\
$J_t^\top Q$, correct indexing & 6/9 & 3 \\
Historical $J_{t+1}^\top Q$ & 7/9 & 2 \\
\bottomrule
\end{tabular}
\end{center}

The two remaining historical-analysis differences are both adjacent grid
points at $N=1000$. They are much smaller than the discrepancies produced by
the corrected analysis. At $N=1000,\alpha=1$, for example, the historical
mean MLE is $-0.0062$ at the paper's $g=0.791$ and $+0.0118$ at the next
gain, $g=0.910$. The reported grid point is therefore sensitive to the exact
set of ten random realizations.

\section{Finding 2: the corrected analysis deviates from Figure 2}

At four informative paper boundaries, changing from the corrected analysis to
the correctly indexed transposed analysis changes the mean MLE as follows:

\begin{center}
\begin{tabular}{rrrrr}
\toprule
$N$ & $\alpha$ & $g$ & Correct $J_tQ$ & $J_t^\top Q$ \\
\midrule
1,000 & 1.0 & 0.791 & $+0.1470$ & $-0.0002$ \\
1,000 & 1.5 & 0.791 & $+0.0958$ & $+0.0167$ \\
3,000 & 1.0 & 0.450 & $+0.0985$ & $+0.0131$ \\
10,000 & 1.0 & 0.295 & $+0.0496$ & $-0.0008$ \\
\bottomrule
\end{tabular}
\end{center}

The transpose changes the mean MLE by roughly $0.05$ to $0.15$ in these
heavy-tailed conditions. By comparison, shifting the Jacobian index from the
correct $J_t$ to the historical $J_{t+1}$ changes the same means by only
about $0.002$ to $0.006$. The indexing error is nevertheless sufficient to
move $N=10000,\alpha=1$ from $g^*=0.339$ to the paper's $g^*=0.295$,
because the mean lies extremely close to zero.

\section{Why the largest discrepancies occur for small heavy-tailed networks}

The corrected and historical analyses disagree most for $N=1000$ with
$\alpha=1.0$ or $1.5$. Three effects reinforce one another there.

\subsection{Nonlinear and heterogeneous Jacobian gating}

Writing $J_t=D_tW$ makes the orientation error explicit:

\[
J_tQ=D_tWQ,
\qquad
J_t^\top Q=W^\top D_tQ.
\]

The two updates gate opposite sides of the recurrent matrix. At the paper's
reported gains, the $N=1000$ heavy-tailed networks contain substantially more
saturated activity, so $D_t$ is heterogeneous and far from the identity:

\begin{center}
\begin{tabular}{rrrr}
\toprule
$N$ & $\alpha$ & Mean $1-h^2$ & Fraction with $|h|\geq0.99$ \\
\midrule
1,000 & 1.0 & 0.738 & 7.43\% \\
3,000 & 1.0 & 0.932 & 1.66\% \\
10,000 & 1.0 & 0.987 & 0.29\% \\
1,000 & 1.5 & 0.810 & 1.72\% \\
3,000 & 1.5 & 0.955 & 0.31\% \\
10,000 & 1.5 & 0.984 & 0.12\% \\
\bottomrule
\end{tabular}
\end{center}

At larger $N$, the heavy-tail transition occurs at a lower gain while the
activity remains closer to the quiescent state. Consequently $D_t\approx I$,
and the orientation error has less effect on the detected crossing. Gaussian
networks also lack the rare, dominant connections that make heavy-tailed
Jacobians strongly heterogeneous and non-normal.

\subsection{Weak self-averaging}

For $\alpha<2$, rare extreme weights dominate individual network
realizations. Transition locations therefore concentrate slowly with network
size. In the corrected sweep, the standard deviation of the per-trial crossing
index was 1.64, 1.55, and 0.46 grid steps for $\alpha=1$ at
$N=1000,3000,10000$, respectively. For $\alpha=1.5$, the corresponding
values were 0.87, 0.75, and 0.50. Gaussian values were only 0.40--0.49 grid
steps.

\subsection{Coarse thresholding and different random streams}

The paper reports the first nonnegative point on a 50-point gain grid rather
than an interpolated root. Small changes around zero therefore move the
reported result by a full 15\% grid step. The original study used a sequential
SciPy/global RNG stream beginning with \texttt{seed + trial}
(\href{original/run_figs/fig2_parallel.py\#L39-L64}{\texttt{fig2\_parallel.py}}); this
reimplementation uses independently keyed condition streams. Both sample the
same target distributions, but they do not generate the same ten matrices.

This distinction matters most for the two $N=1000$ heavy-tail cases. Under
the historical re-analysis, a bootstrap over the current ten
$N=1000,\alpha=1$ trial curves selected the paper's $g=0.791$ crossing about
20\% of the time. The archived submission script also specifies only
$N=1000$ and trial 0
(\href{original/submit_jobs/submit_fig2_parallel.sh\#L5-L9}{\texttt{submit\_fig2\_parallel.sh}}),
not the three sizes and ten trials reported in the paper, so it does not retain
the exact execution provenance of Figure 2.

\section{Interpretation}

The completed flagship sweep supports two distinct reproducibility claims:

\begin{itemize}
\item \textbf{Historical numerical reproduction:} yes. When the historical transpose
and indexing behavior are restored, the reported Figure 2 transition points
are recovered to within one sampled gain everywhere, with exact agreement in
7 of 9 conditions.
\item \textbf{Reproduction under the algorithm stated in the paper:} no, not
quantitatively. The corrected $J_tQ$ analysis yields materially earlier
transitions for heavy-tailed networks, most prominently at
$N=1000,\alpha=1$, $N=1000,\alpha=1.5$, and $N=3000,\alpha=1$.
\end{itemize}

The corrected sweep still reproduces the qualitative scientific pattern:
heavier tails yield a slower, broader approach to chaos, and increasing
network size shifts the transition toward lower gain. What changes is the
quantitative location and magnitude of that transition. Accordingly, results
from \texttt{runs/flagship/} should be described as a \textbf{corrected-analysis
replication}, while the historical re-analysis strongly suggests how the
paper's published numerical values arose.

\end{document}
